\section{Legal Language Models and Benchmarks (2017--2023)}
\label{sec:lm}

\subsection{From manually designed rules to learned legal representations}

The transition from knowledge-based systems to neural language models changed how computers
processed legal text. Earlier systems depended heavily on explicit rules, predefined categories, or
manually designed features. Neural models instead learned statistical representations from examples,
allowing them to identify linguistic patterns that were difficult to encode manually.

This transition did not make legal knowledge representation obsolete. Rather, it changed where much
of the representation was learned. A neural model could learn that particular expressions frequently
occur in contracts, judgments, or legislation without a developer specifying every linguistic
pattern in advance. Such representations could support tasks including document classification,
legal information extraction, question answering, and case retrieval.

The difficulty was that legal language differs substantially from everyday language. Statutes use
carefully defined terms; contracts contain specialized expressions; judgments distinguish between
the facts, arguments, reasoning, and decisions of a case. A model trained predominantly on general
text may understand the grammar of a legal sentence while failing to represent the significance of
its terminology.

Pretrained transformer models offered a way to address this problem. Rather than training a language
model from scratch for every legal task, researchers could adapt a pretrained model to legal
documents and then fine-tune it for a specific application. This approach reduced the need to build a
separate linguistic system for every task, although the resulting performance still depended on the
training data, model architecture, and evaluation setting.

\subsection{Domain-specific models: LEGAL-BERT}

A significant development was the introduction of LEGAL-BERT by \citet{chalkidis2020}. The study
investigated different strategies for adapting BERT to legal NLP tasks, including using the original
pretrained model, continuing pretraining on domain-specific text, and pretraining from scratch on
legal corpora. The authors released a family of legal language models intended to support research
in computational law and legal technology.

The underlying principle was straightforward: if a model must process legal documents, its training
should expose it to the language and terminology found in those documents. Domain adaptation can
help a model learn distinctions that are less common in general-purpose text. It also provides a more
systematic starting point for multiple downstream tasks than developing a new model for each one.

However, LEGAL-BERT should not be interpreted as a complete legal reasoning system. It provides
learned language representations that can support particular NLP tasks; it does not automatically
establish which legal authority governs a dispute, whether a provision remains applicable, or whether
a generated legal conclusion is supported by the cited text. These questions require additional
retrieval, legal context, and verification mechanisms.

This distinction is important for understanding the subsequent development of legal RAG. A
domain-specific language model improves how a system processes legal language, but it does not
necessarily provide access to current authoritative documents at inference time. Model knowledge and
retrieved legal evidence serve different purposes.

\subsection{LexGLUE: evaluating legal language understanding}

As legal language models became more common, another problem became apparent: how could researchers
determine whether a model was genuinely effective across legal tasks rather than simply successful on
a single dataset?

\citet{chalkidis2022} addressed this problem with LexGLUE, the Legal General Language Understanding
Evaluation benchmark. LexGLUE brings together multiple datasets for evaluating legal language
understanding in a standardized setting. Its purpose is to compare general-purpose and legal-oriented
language models across different tasks rather than relying on results from a single application. The
authors reported that legal-oriented models consistently improved performance over general-purpose
models across multiple tasks in their evaluation.

The importance of LexGLUE lies in its emphasis on comparative evaluation. A model that performs well
on one legal classification task may not perform equally well on another task requiring different
linguistic or legal distinctions. A shared benchmark makes these differences easier to investigate
and provides a common basis for future experiments.

Nevertheless, benchmark performance has a defined scope. LexGLUE primarily evaluates legal language
understanding on its constituent datasets. Strong results do not automatically demonstrate that a
system can retrieve the correct authority for a new legal question, preserve the hierarchy of a
statute, generate a legally supported explanation, or verify citations. These are related
capabilities, but they are not interchangeable.

\subsection{LegalBench: evaluating legal reasoning}

LegalBench extended the evaluation discussion toward the kinds of reasoning tasks that legal
professionals consider important. \citet{guha2023} introduced LegalBench, a collaboratively
constructed benchmark containing 162 tasks covering six types of legal reasoning. The tasks were
developed through collaboration between legal and computational researchers, with legal professionals
contributing to their design. The study evaluated 20 open-source and commercial LLMs and examined how
the tasks correspond to established ways of describing legal reasoning.

The benchmark addressed a central limitation of evaluating legal AI through generic language
benchmarks alone: a model may perform well on ordinary language understanding while struggling with
tasks that require legally meaningful distinctions. LegalBench provides a more direct way to examine
whether models can carry out a range of tasks associated with legal analysis.

At the same time, the existence of a legal reasoning benchmark does not establish that a model is
suitable for unsupervised legal practice. Performance depends on the task definitions, examples,
evaluation metrics, and tested models. A benchmark can reveal strengths and weaknesses, but it cannot
cover every jurisdiction, factual scenario, procedural context, or change in the law.

LegalBench and LexGLUE therefore serve complementary purposes. LexGLUE provides a standardized
evaluation framework for legal language understanding, while LegalBench focuses on a broader
collection of tasks designed to examine legal reasoning. Neither benchmark, by itself, resolves the
problem of grounding a generated answer in the correct, current, and authoritative legal source.

\subsection{What changed, and what remained unresolved?}

The developments between 2017 and 2023 established two important capabilities. First, domain-adapted
pretrained models made legal language processing more flexible and reusable. Second, shared
benchmarks provided more systematic ways to compare models and investigate their limitations.

However, three challenges remained. First, learned representations do not guarantee access to
authoritative evidence. A model may have learned patterns from legal text but cannot be assumed to
know whether a statute has subsequently been amended or which authority applies to a particular
dispute. Second, benchmark success is not equivalent to legal reliability. Classification accuracy or
performance on a reasoning task does not establish that every claim in a generated response is
factually supported or that its citations are correct. Third, legal context remains important. A
model may process a paragraph accurately while overlooking the relationship between that paragraph and
its parent provision, an exception, an amendment, or a related judicial interpretation.

These limitations help explain why the next phase of LegalAI increasingly incorporated retrieval into
answer generation. Rather than depending exclusively on information encoded in model parameters,
retrieval-augmented systems could search an external legal corpus and provide relevant documents to
the generator. The aim was to combine the language capabilities of neural models with access to
explicit evidence. Table~\ref{tab:lmdev} summarizes the main developments of this period.

\begin{table*}[tp]
\centering
\caption{Major developments in legal language models and evaluation (2017--2023).}
\label{tab:lmdev}
\footnotesize
\setlength{\tabcolsep}{5pt}
\renewcommand{\arraystretch}{1.3}
\begin{tabularx}{\textwidth}{>{\raggedright\arraybackslash}p{2.8cm} >{\raggedright\arraybackslash}X >{\raggedright\arraybackslash}X >{\raggedright\arraybackslash}X}
\toprule
\textbf{Development} & \textbf{Main contribution} & \textbf{What it enabled} & \textbf{What it did not establish} \\
\midrule
Domain-adapted neural models & Learned representations from legal language and task-specific examples. & More flexible processing of legal text than systems dependent entirely on manually designed features. & Correct legal interpretation or access to current authoritative sources. \\
LEGAL-BERT \citep{chalkidis2020} & Investigated BERT adaptation strategies and released a family of legal language models. & Reusable pretrained representations for legal NLP research. & End-to-end legal reasoning, citation verification, or guaranteed factual correctness. \\
LexGLUE \citep{chalkidis2022} & Standardized evaluation across multiple legal language understanding datasets. & More consistent comparison of general-purpose and legal-oriented models across tasks. & Reliable retrieval and generation of evidence-grounded legal answers. \\
LegalBench \citep{guha2023} & Introduced 162 collaborative legal reasoning tasks spanning six types of reasoning. & More direct investigation of LLM capabilities on tasks designed to reflect legal reasoning. & Universal competence across jurisdictions or dependable performance in real legal practice. \\
\bottomrule
\end{tabularx}
\end{table*}

The central lesson of this period is that better legal language understanding and reliable legal
assistance are related, but different, goals. Domain-specific models improved how systems processed
legal text, while benchmarks made their capabilities easier to measure. Yet neither development alone
guaranteed that a system could find the right authority, retain its legal context, and justify its
answer with verifiable evidence.

That unresolved problem motivates the next stage of the survey: the emergence of legal
retrieval-augmented generation, where retrieval becomes an explicit part of the answer-generation
process.
